\section{Discussion}
\label{sec:discussion}

\paragraph{What makes an agent strong.} Every layer returns the same answer: strong agents win by writing queries that retrieve, and nearly everything else is downstream. Accuracy is dominated by retrieval recall rather than context budget; within the same question, the agent that retrieves the gold evidence usually answers correctly, while the one that misses it almost never does. Recall also absorbs most of the apparent between-agent capability gap (§\ref{sec:results-synthesis}). The behavioral signature of strength is discipline, not effort: strong agents essentially never re-ask the same content, pivot as often as weak agents but \emph{land} those pivots, often batch several angles per turn, and stop near evidence saturation. High search volume is therefore a symptom of a failing trajectory, not a lever that rescues it.

\paragraph{What transfers to weaker agents, and how.} Similar accuracies can hide disjoint failure profiles: Kimi K2.6 is utilization-bound while Deepseek V4 Pro is retrieval-bound at the same accuracy, and Tongyi-DR's apparent search problem is largely a missing stopping rule. Transfer should therefore be prescribed by diagnosis, not leaderboard position. The largest transferable surface is the query side. On the shared retriever the $26$-point spread in gold recall is query-side by construction (§\ref{sec:results-failure}), so retrieval-gap errors are not benchmark-impossible but query-fixable in principle. Possible mechanisms include reformulation and deduplication guards \citep{huang2009analyzing}, prompt-level strategy transfer, or distillation on stronger agents' queries, though our numbers are upper bounds rather than realized gains. Other levers are scaffold-side and cheap: answer normalization for format-or-subtle errors, schema repair for rejected searches, snippet-stream management, and snippet-first visit gates \citep{Chen2026AgentIRReasoning-AwareRetrieval}.

\paragraph{What remains at the frontier.} For the strongest agents the diagnosis shifts. GLM 5.1 sits closest to the oracle reader ceiling with a negligible wasted tail, so reformulation triggers and stopping rules have little left to reclaim. Its residual errors concentrate in mis-anchoring (Q2 of Table~\ref{tab:query-case}) and in the format-or-subtle share, where the prescriptions are verification and answer normalization rather than more search. This suggests a different frontier problem: not how to spend more retrieval budget, but how to decide when an apparently coherent chain is anchored to the wrong entity, and how to normalize answers without weakening precision. Complementarity also matters: agents fail on largely disjoint questions, so no single agent's error set marks the field's ceiling. The oracle reader places that ceiling at $87.5\%$ (§\ref{sec:results-failure}), leaving a residual that gold evidence in hand does not fix.

\paragraph{Relation to human search behavior.} The patterns also clarify where current agents differ from human searchers. Berrypicking \citep{bates1989berrypicking} predicts that information needs evolve as evidence accumulates; our agents often do the opposite, fixing an initial framing and then issuing paraphrase searches long after new evidence has stopped arriving \citep{huang2009analyzing, rieh2006analysis}. Humans also stop under satisficing, fatigue, and the economics of search cost \citep{zach2005enough, azzopardi2014economic}; agents usually stop at hard caps or after emitting an answer. The missing piece is not merely a shorter budget, but an internal sense of marginal utility: whether a new query is likely to change the answer, whether a snippet adds new evidence, and whether enough independent support has accumulated. In that sense, evidence-driven stopping and query reformulation are not just engineering tricks; they are the agentic analogues of human search judgment.
